\section{Geometry}
\subsection{Precision Handling}
\textbf{Note:} Never use \texttt{a == b} for floats. Use \texttt{abs(a - b) < EPS}.
\begin{minted}{c++}
const ld EPS = 1e-9;
bool equal(ld a, ld b) { return abs(a - b) < EPS; }
bool lessThan(ld a, ld b) { return a < b - EPS; }
bool greaterThan(ld a, ld b) { return a > b + EPS; }
\end{minted}
\subsection{Geometry Formulas}
\textbf{General Triangle} (sides $a, b, c$; $s = \frac{a+b+c}{2}$): \\
\textbullet{} \textbf{Area:} $A = \frac{1}{2}bh = \frac{1}{2}ab\sin C = \sqrt{s(s-a)(s-b)(s-c)}$ \\
\textbullet{} \textbf{Radii:} Inradius $r = \frac{A}{s}$, \quad Circumradius $R = \frac{abc}{4A}$ \\
\textbullet{} \textbf{Heights ($h_a$):} $\perp$ to base $a$: $h_a = \frac{2A}{a}$; \quad $\sum \frac{1}{h_i} = \frac{1}{r}$ \\
\textbullet{} \textbf{Excircle ($r_a$):} Outside touching $a$: $r_a = \frac{A}{s-a}$; \quad $\sum \frac{1}{r_i} = \frac{1}{r}$ \\
\textbullet{} \textbf{Sine/Cosine:} $\frac{a}{\sin A} = \frac{b}{\sin B} = 2R$; \ $\cos C = \frac{a^2+b^2-c^2}{2ab}$ \\
\textbullet{} \textbf{Medians ($m_a$):} To midpoint of $a$: $m_a = \frac{1}{2}\sqrt{2b^2+2c^2-a^2}$ \\
\textbullet{} \textbf{Area from 3 Medians:} $A = \frac{4}{3} \times \text{Heron}(m_a, m_b, m_c)$ \\
\textbullet{} \textbf{Angle Bisector:} $t_a = \sqrt{bc - mn} = \frac{2bc\cos(A/2)}{b+c}$ ($\frac{m}{n} = \frac{c}{b}, \ m+n=a$) \\
\textbullet{} \textbf{Stewart's Thm:} Cevian $d$ on $a = m+n$: $b^2 m + c^2 n = a(d^2 + mn)$ \\
\textbf{Special Triangles}: \\
\textbullet{} \textbf{Equilateral}: $h = \frac{\sqrt{3}}{2}a, \ A = \frac{\sqrt{3}}{4}a^2, \ r = \frac{a}{2\sqrt{3}}, \ R = 2r$ \\
\textbullet{} \textbf{Isosceles} (base $b$, equal $a$): $h = \sqrt{a^2 - \frac{b^2}{4}}, \ A = \frac{b}{4}\sqrt{4a^2-b^2}$ \\
\textbullet{} \textbf{Right-Angled}: $A = \frac{1}{2}ab, \ r = \frac{a+b-c}{2}, \ R = \frac{c}{2}, \ h = \frac{ab}{c}$ \\

Circles:
Arc Length: $L = \theta r$
Sector Area: $Area = \frac{1}{2} r^2 \theta$
Segment Area: $Area = \frac{1}{2} r^2 (\theta - \sin\theta)$

\textbf{Rhombus:} Area = $\frac{1}{2} \times d_1 \times d_2$ \\
\textbf{Trapezium:} Area = $\frac{1}{2} \times (b_1 + b_2) \times h$ \\
\textbf{Trapezium ($a \parallel c$, legs $b,d$):} $z = \frac{b^2-d^2+(a-c)^2}{2(a-c)}$, $h = \sqrt{b^2-z^2}$, $A = \frac{a+c}{2}h$ \\
\textbf{Regular Polygon (n sides):} Perimeter = $n \times a$, Area = $\frac{1}{4}na^2 \cot(\frac{\pi}{n})$ \\
\textbf{Circle:} $C = 2\pi r$, $A = \pi r^2$ \\
\textbf{Ellipse:} $A = \pi a b$ \\
\textbf{Polygon:} Interior sum = $(n-2) \times 180^\circ$

3D Shapes:
\textbf{Sphere:} $SA = 4\pi r^2$, $V = \frac{4}{3}\pi r^3$ \\
\textbf{Cylinder:} $SA = 2\pi r(r+h)$, $V = \pi r^2 h$ \\
\textbf{Cone:} $SA = \pi r(r+l)$, $V = \frac{1}{3}\pi r^2 h$ \\
\textbf{Conical Frustum:} $V = \frac{\pi}{3} h (r_1^2 + r_2^2 + r_1 r_2)$ \\
\textbf{Torus:} $SA = 4\pi^2 Rr$, $V = 2\pi^2 Rr^2$ \\
\textbf{Cube:} $V = |x_2-x_1| \cdot |y_2-y_1| \cdot |z_2-z_1|$ \\

\subsection{Geometry Theorems}
\textbf{Pick's Theorem}: For a simple polygon with integer coordinates:
$A = I + \frac{B}{2} - 1$, where $I$ = interior lattice points, $B$ = boundary lattice points.

\textbf{Brahmagupta's Formula}: Area of a cyclic quadrilateral with sides $a,b,c,d$:
$S = \sqrt{(s-a)(s-b)(s-c)(s-d)}$ where $s = \frac{a+b+c+d}{2}$.

\textbf{Fermat Point}: Point minimizing sum of distances to triangle vertices.
If no angle $> 120^\circ$, sum is $\sqrt{(a^2+b^2+c^2 + 4\sqrt{3}\times Area)/2}$.

\textbf{Simpson's Rule}: $\int_a^b f(x) dx \approx \frac{h}{3}[f(x_0) + f(x_n) + 4\sum f(x_{2i+1}) + 2\sum f(x_{2i})]$

\textbf{Pythagorean Triples}: $a = 2st, b = s^2-t^2, c = s^2+t^2$ for $s>t>0$, $\gcd(s,t)=1$, $s \not\equiv t \pmod 2$.

\subsection{Closest Pair of Points}
\textbf{Returns:} \texttt{Res} struct containing the minimum distance squared (\texttt{d}) and the two points (\texttt{a}, \texttt{b}).
\begin{minted}{c++}
struct Res { int64_t d; P a, b; };
int64_t sq(int64_t x) { return x * x; }
int64_t dist_sq(P &a, P &b) {
  return sq(a.x - b.x) + sq(a.y - b.y);
}
Res closest_pair(vector<P>& p, vector<P>& py,
vector<P>& tmp, int l, int r) {
  int n = r - l + 1;
  if (n <= 1) return {LLONG_MAX, {}, {}};
  if (n == 2) {
    if (py[l].y > py[l + 1].y)
      swap(py[l], py[l + 1]);
    return {dist_sq(p[l], p[l + 1]), p[l], p[l + 1]};
  }
  int m = (l + r) >> 1, midx = p[m].x;
  Res res_l = closest_pair(p, py, tmp, l, m);
  Res res_r = closest_pair(p, py, tmp, m + 1, r);
  Res res = (res_l.d < res_r.d ? res_l : res_r);
  for (int i = l, j = m + 1, k = l; k <= r; k++) {
    if (j > r || (i <= m && py[i].y < py[j].y))
      tmp[k] = py[i++];
    else tmp[k] = py[j++];
  }
  for (int i = l; i <= r; i++) py[i] = tmp[i];
  vector<P> stripe;
  for (int i = l; i <= r; i++)
    if (sq(py[i].x - midx) < res.d)
      stripe.push_back(py[i]);
  for (int i = 0; i < (int)stripe.size(); i++) {
    for (int j = i + 1; j < (int)stripe.size(); j++) {
      if (sq(stripe[i].y - stripe[j].y) >= res.d)
        break;
      int64_t d = dist_sq(stripe[i], stripe[j]);
      if (d < res.d)
        res = {d, stripe[i], stripe[j]};
    }
  }
  return res;
}
\end{minted}
\subsection{Rectangle Distance}
Distance between the two incircle touch points of the right triangles formed by the diagonal of an $L \times W$ rectangle.
\begin{minted}{c++}
ld L, W; if (!(cin >> L >> W)) return 0;
ld h = sqrtl(L*L + W*W), P = L + W + h;
ld x1 = L*(L + h)/P, y1 = L*W/P;
ld x2 = y1, y2 = W*(W + h)/P;
ld dx = x1 - x2, dy = y1 - y2;
ld D = sqrtl(dx*dx + dy*dy);
cout << fixed << setprecision(6) << (double)D << "\n";
\end{minted}
\subsection{Polygon Operations}
\begin{minted}{c++}
const ld EPS = 1e-9;
const ld PI = acos(-1.0);
// sign: -1/0/1 for negative/zero/positive
int sgn(ld x) { return (x > EPS) - (x < -EPS); }
// 2D point with int64 coords: cross(*), dot(^),
// norm(sq), abs, rotation
struct P {
  using T = int64_t;
  T x, y;
  void read() { cin >> x >> y; }
  P operator - (const P &b) const { return P{x - b.x, y - b.y}; }
  P operator + (const P &b) const { return P{x + b.x, y + b.y}; }
  template<class U> P operator * (const U &c) const {
    return P{(T)(x * c), (T)(y * c)}; }
  template<class U> P operator / (const U &c) const {
    return P{(T)(x / c), (T)(y / c)}; }
  bool operator < (const P &b) const { return x == b.x ? y < b.y : x < b.x; }
  bool operator == (const P &b) const { return x == b.x && y == b.y; }
  T operator * (const P &b) const { return x * b.y - y * b.x; }
  T operator ^ (const P &b) const { return x * b.x + y * b.y; }
  T dot(P &a, P &b) const { return (a - *this) ^ (b - *this); }
  T cross(P &a, P &b) { return (a - *this) * (b - *this); }
  ld sq() const { return x * x + y * y; }
  ld abs() const { return sqrtl(sq()); }
  P rot() const { return P(-y, x); } // 90 CCW
  P rot(ld a) const {
    return P(x*cos(a) - y*sin(a),
             x*sin(a) + y*cos(a)); }
};
// true if a, b, c lie on same line
bool collinear(P &a, P &b, P &c) {
  return (b - a) * (c - a) == 0;
}
// true if r lies on closed segment pq
bool on_segment(P &p, P &q, P &r) {
  return collinear(p, q, r) &&
    r.x <= max(p.x, q.x) && r.x >= min(p.x, q.x) &&
    r.y <= max(p.y, q.y) && r.y >= min(p.y, q.y);
}
// Euclidean distance between a and b
ld dist_euclid(P a, P b) { return (a - b).abs(); }
// project c onto line ab (returns foot point)
P project(P a, P b, P c) {
  return a + (b - a) * (((c - a) ^ (b - a)) /
         (b - a).sq());
}
// reflect c across line ab
P reflect(P a, P b, P c) {
  return project(a, b, c) * 2 - c;
}
// perpendicular dist from c to line ab
ld dist_pt_line(P a, P b, P c) {
  return fabsl((b - a) * (c - a)) / (b - a).abs();
}
// shortest dist from c to segment ab
ld dist_pt_seg(P a, P b, P c) {
  if (((c - a) ^ (b - a)) <= 0) return (c - a).abs();
  if (((c - b) ^ (a - b)) <= 0) return (c - b).abs();
  return dist_pt_line(a, b, c);
}
// true if p lies on segment ab (EPS tolerance)
bool on_seg(P a, P b, P p) {
  return !sgn((b - a) * (p - a)) &&
         ((p - a) ^ (p - b)) <= EPS;
}
// line-line intersection (ab x cd), stores in res,
// returns false if parallel
bool line_int(P a, P b, P c, P d, P &res) {
  ld cp = (b - a) * (d - c);
  if (!sgn(cp)) return false;
  res = a + (b - a) * (((c - a) * (d - c)) / cp);
  return true;
}
// true if segment ab intersects segment cd
bool seg_int(P a, P b, P c, P d) {
  auto cp = [&](P x, P y, P z) {
    return sgn((y - x) * (z - x)); };
  if (on_seg(a, b, c) || on_seg(a, b, d) ||
      on_seg(c, d, a) || on_seg(c, d, b))
    return true;
  return cp(a, b, c) != cp(a, b, d) &&
         cp(c, d, a) != cp(c, d, b);
}
// area of polygon (shoelace), vertices in order (strict CCW or CW)
ld polygon_area(vector<P> &p) {
  int64_t area = 0, n = p.size();
  for (int i = 0; i < n; ++i)
    area += p[i] * p[(i + 1) % n];
  return (ld)abs(area) / 2.0L;
}
// convex hull vertices (CCW, monotone chain)
vector<P> convexHull(vector<P> p) {
  int n = p.size(), k = 0;
  if (n <= 2) return p;
  vector<P> h(2 * n);
  sort(p.begin(), p.end());
  for (int i = 0; i < n; h[k++] = p[i++])
    while (k >= 2 && h[k - 2].cross(h[k - 1],
           p[i]) <= 0) k--;
  for (int i = n - 2, t = k + 1; i >= 0;
       h[k++] = p[i--])
    while (k >= t && h[k - 2].cross(h[k - 1],
           p[i]) <= 0) k--;
  h.resize(k - 1);
  return h;
}
// point-in-polygon: 1=inside, 2=on boundary, 0=out
int is_point_in_polygon(const vector<P>& p, P q) {
  int c = 0;
  for (int i = 0, n = p.size(); i < n; ++i) {
    P a = p[i], b = p[(i + 1) % n];
    if ((b - a) * (q - a) == 0 &&
        min(a.x, b.x) <= q.x &&
        q.x <= max(a.x, b.x) &&
        min(a.y, b.y) <= q.y &&
        q.y <= max(a.y, b.y)) return 2;
    if ((a.y > q.y) != (b.y > q.y) &&
        ((b - a) * (q - a) > 0) == (a.y < b.y))
      c = !c;
  }
  return c;
}
\end{minted}
\subsection{Area Enclosed by Three Circles}
Black area between three mutually tangent circles of radii $r_1, r_2, r_3$ (triangle area minus sector areas).
\begin{minted}{c++}
ld r1, r2, r3; cin >> r1 >> r2 >> r3;
ld a = r2 + r3, b = r3 + r1, c = r1 + r2;
ld s = (a + b + c) / 2.0L;
ld val = s*(s - a)*(s - b)*(s - c);
ld T = val > 0 ? sqrtl(val) : 0;
auto angle = [&](ld opp, ld x, ld y) {
  ld cosv = (x*x + y*y - opp*opp) / (2.0L * x * y);
  return acosl(min((ld)1.0, max((ld)-1.0, cosv)));
};
ld A = angle(a, b, c), B = angle(b, c, a);
ld C = angle(c, a, b);
ld sectorSum = 0.5L*(r1*r1*A + r2*r2*B + r3*r3*C);
ld blackArea = T - sectorSum;
cout << fixed << setprecision(6) <<
(double)blackArea << "\n";
\end{minted}
\subsection{Half-Plane Intersection}
Finds the convex polygon (or empty region) satisfying all half-planes $a_i x + b_i y \le c_i$.
Lines are sorted by angle, then a deque maintains the intersection.
\texttt{halfPlaneIntersection(lines)} returns vertices of the resulting polygon. $O(N \log N)$.
\begin{minted}{c++}
struct Line {
  double a, b, c, ang; // ax + by <= c
  Line(double a = 0, double b = 0, double c = 0)
    : a(a), b(b), c(c), ang(atan2(b, a)) {}
  bool operator < (const Line &o) const { return ang < o.ang; }
};
pair<double,double> intersect(Line &p, Line &q) {
  double d = p.a * q.b - p.b * q.a;
  return {(p.c * q.b - q.c * p.b) / d,
          (p.a * q.c - q.a * p.c) / d};
}
bool ok(Line &a, Line &b, Line &c) {
  auto [x, y] = intersect(a, b);
  return c.a * x + c.b * y < c.c + 1e-9;
}
vector<pair<double,double>> halfPlaneIntersection(vector<Line> v) {
  sort(v.begin(), v.end());
  vector<Line> filtered;
  for (int i = 0; i < (int)v.size(); i++) {
    if (i > 0 && abs(v[i].ang - v[i-1].ang) < 1e-9) {
      if (v[i].c < filtered.back().c) filtered.back() = v[i];
    } else filtered.push_back(v[i]);
  }
  v = filtered;
  deque<Line> dq;
  for (auto &l : v) {
    while (dq.size() >= 2 && !ok(dq[dq.size()-2], dq.back(), l)) dq.pop_back();
    while (dq.size() >= 2 && !ok(dq[0], dq[1], l)) dq.pop_front();
    dq.push_back(l);
  }
  while (dq.size() >= 3 && !ok(dq[dq.size()-2], dq.back(), dq[0])) dq.pop_back();
  while (dq.size() >= 3 && !ok(dq[0], dq[1], dq.back())) dq.pop_front();
  vector<pair<double,double>> res;
  for (int i = 0; i < (int)dq.size(); i++)
    res.push_back(intersect(dq[i], dq[(i+1)%dq.size()]));
  return res;
}
\end{minted}

\subsection{Rotating Calipers}
Extends convex hull: $O(N)$ diameter (farthest pair) and minimum width. Requires CCW hull from \texttt{convexHull}.
\begin{minted}{c++}
// Diameter: farthest pair in convex polygon, returns actual distance
ld rotating_diameter(vector<P> &h) {
  int n = h.size();
  if (n <= 1) return 0;
  if (n == 2) return (h[0] - h[1]).abs();
  ld ans = 0;
  for (int i = 0, j = 1; i < n; i++) {
    while ((h[(i+1)%n] - h[i]) * (h[(j+1)%n] - h[j]) > 0) j = (j + 1) % n;
    ans = max({ans, (h[i] - h[j]).sq(), (h[(i+1)%n] - h[j]).sq()});
  }
  return sqrtl(ans);
}
// Width: minimum distance between parallel supporting lines
ld polygon_width(vector<P> &h) {
  int n = h.size();
  if (n <= 2) return 0;
  ld ans = 1e30;
  for (int i = 0, j = 1; i < n; i++) {
    while ((h[(i+1)%n] - h[i]) * (h[(j+1)%n] - h[j]) > 0) j = (j + 1) % n;
    ld d = fabsl((h[(i+1)%n] - h[i]) * (h[j] - h[i])) / (h[(i+1)%n] - h[i]).abs();
    ans = min(ans, d);
  }
  return ans;
}
\end{minted}

\subsection{Point-in-Convex Polygon (O(log N))}
Determines if a point is inside a CCW-ordered convex hull in $O(\log N)$.
\begin{minted}{c++}
// Returns 1 if q is inside or on boundary of convex polygon p (CCW), 0 otherwise
int point_in_convex(P &q, vector<P> &p) {
  int n = p.size();
  auto in_box = [&](P a, P b) {
    return min(a.x, b.x) <= q.x && q.x <= max(a.x, b.x) &&
           min(a.y, b.y) <= q.y && q.y <= max(a.y, b.y);
  };
  int64_t cross_first = (p[1] - p[0]) * (q - p[0]);
  if (cross_first == 0) return in_box(p[0], p[1]);
  if (cross_first < 0) return 0;
  int lo = 1, hi = n - 1;
  while (lo < hi) {
    int mid = (lo + hi + 1) / 2;
    if ((p[mid] - p[0]) * (q - p[0]) >= 0) lo = mid;
    else hi = mid - 1;
  }
  int64_t cross_last = (p[(lo+1)%n] - p[lo]) * (q - p[lo]);
  if (cross_last < 0) return 0;
  if (cross_last == 0) return in_box(p[lo], p[(lo+1)%n]);
  return 1;
}
\end{minted}

\subsection{Minimum Enclosing Circle (Welzl)}
Randomized expected $O(N)$. Finds smallest circle enclosing all points.
\begin{minted}{c++}
struct PointD {
  ld x, y;
  PointD(ld x=0, ld y=0) : x(x), y(y) {}
  PointD(P p) : x(p.x), y(p.y) {}
  PointD operator+(PointD p) const { return {x+p.x, y+p.y}; }
  PointD operator-(PointD p) const { return {x-p.x, y-p.y}; }
  PointD operator*(ld c) const { return {x*c, y*c}; }
  PointD operator/(ld c) const { return {x/c, y/c}; }
  ld cross(PointD p) const { return x*p.y - y*p.x; }
  ld sq() const { return x*x + y*y; }
  ld abs() const { return sqrtl(sq()); }
};
struct Circle { PointD c; ld r; };
Circle welzl(vector<P> &pts) {
  mt19937 rng(chrono::steady_clock::now()
    .time_since_epoch().count());
  vector<PointD> p(pts.begin(), pts.end());
  shuffle(p.begin(), p.end(), rng);
  auto circum = [](PointD a, PointD b, PointD c) {
    PointD ab = b-a, ac = c-a;
    ld d = 2 * ab.cross(ac);
    PointD o(
      a.x + (ab.sq()*ac.y - ac.sq()*ab.y) / d,
      a.y + (ac.sq()*ab.x - ab.sq()*ac.x) / d);
    return Circle{o, (o - a).abs()};
  };
  function<Circle(int, vector<PointD>&)> f =
    [&](int i, vector<PointD> &b) -> Circle {
    if (i == (int)p.size() || b.size() == 3) {
      if (b.empty()) return {{0,0}, 0};
      if (b.size() == 1) return {b[0], 0};
      if (b.size() == 2) return {(b[0]+b[1])/2, (b[0]-b[1]).abs()/2};
      return circum(b[0], b[1], b[2]);
    }
    Circle c = f(i+1, b);
    if ((c.c - p[i]).abs() > c.r + 1e-9) {
      b.push_back(p[i]);
      c = f(i+1, b);
      b.pop_back();
    }
    return c;
  };
  vector<PointD> b;
  return f(0, b);
}
// welzl(p).r is the radius of the minimum enclosing circle
\end{minted}

\subsection{Minkowski Sum of Convex Polygons}
$P \oplus Q = \{p+q : p \in P, q \in Q\}$ in $O(N+M)$. Result is convex.
\textbf{Applications:} convex polygon collision detection, distance computation.
\begin{minted}{c++}
vector<P> minkowski_sum(vector<P> &a, vector<P> &b) {
  auto reorder = [](vector<P> &p) {
    int n = p.size(), idx = 0;
    for (int i = 1; i < n; i++)
      if (p[i].y < p[idx].y ||
         (p[i].y == p[idx].y && p[i].x < p[idx].x))
        idx = i;
    rotate(p.begin(), p.begin()+idx, p.end());
  };
  if (a.empty() || b.empty()) return {};
  reorder(a); reorder(b);
  a.push_back(a[0]); b.push_back(b[0]);
  int n = a.size(), m = b.size();
  vector<P> res;
  for (int i = 0, j = 0; i < n || j < m; ) {
    res.push_back(a[i % n] + b[j % m]);
    if (i >= n) j++;
    else if (j >= m) i++;
    else {
      auto ca = a[(i+1)%n] - a[i%n];
      auto cb = b[(j+1)%m] - b[j%m];
      if (ca * cb >= 0) i++;
      if (ca * cb <= 0) j++;
    }
  }
  return res;
}
\end{minted}

\subsection{Circle Operations}
Circumcircle, circle-line, circle-circle intersection. Uses \texttt{PointD} from Welzl subsection above.
\begin{minted}{c++}
// Circumcircle of three points, returns {center, radius}
pair<PointD, ld> circumcircle(P a, P b, P c) {
  PointD A(a), B(b), C(c);
  PointD ab = B-A, ac = C-A;
  ld d = 2 * ab.cross(ac);
  PointD o(A.x + (ab.sq()*ac.y - ac.sq()*ab.y)/d,
           A.y + (ac.sq()*ab.x - ab.sq()*ac.x)/d);
  return {o, (o - A).abs()};
}
// Circle-Line: returns number of intersection points
int circle_line(PointD c, ld r, P a, P b, vector<PointD> &res) {
  PointD A(a), B(b), ab = B - A;
  PointD p = A + ab * (((c-A).x * ab.x + (c-A).y * ab.y) / ab.sq());
  ld d = (p - c).abs();
  if (d > r + 1e-9) return 0;
  if (fabsl(d - r) < 1e-9) { res = {p}; return 1; }
  ld t = sqrtl(r*r - d*d);
  PointD dir = ab / ab.abs() * t;
  res = {p - dir, p + dir};
  return 2;
}
// Circle-Circle: returns number of intersection points
int circle_circle(PointD c1, ld r1, PointD c2, ld r2, vector<PointD> &res) {
  ld d = (c1 - c2).abs();
  if (d > r1 + r2 + 1e-9 || d < fabsl(r1-r2) - 1e-9) return 0;
  if (d < 1e-9 && fabsl(r1-r2) < 1e-9) return -1;
  ld a = (r1*r1 - r2*r2 + d*d) / (2*d);
  ld h = sqrtl(r1*r1 - a*a);
  PointD m = c1 + (c2-c1) * (a / d);
  PointD perp(-(c2-c1).y * h / d, (c2-c1).x * h / d);
  res = {m + perp, m - perp};
  return 2;
}
\end{minted}

\subsection{Intersection Area of Two Circles}
The overlap region of two intersecting circles forms a symmetric lens bounded by a common chord.
It equals the sum of two circular segments:
\[A = \frac{1}{2}r_1^2(\alpha - \sin\alpha) + \frac{1}{2}r_2^2(\beta - \sin\beta)\]
Central angles in radians:
\[\alpha = 2\cos^{-1}\left(\frac{r_1^2 + d^2 - r_2^2}{2r_1 d}\right),\quad \beta = 2\cos^{-1}\left(\frac{r_2^2 + d^2 - r_1^2}{2r_2 d}\right)\]
If circles do not intersect or one contains the other, the area is 0.
\begin{minted}{c++}
ld circle_intersection_area(ld r1, ld r2, ld d) {
  if (d >= r1 + r2 - EPS) return 0;
  if (d <= fabsl(r1 - r2) + EPS) return 0;
  ld alpha = 2.0L * acosl((r1*r1 + d*d - r2*r2) / (2.0L * r1 * d));
  ld beta  = 2.0L * acosl((r2*r2 + d*d - r1*r1) / (2.0L * r2 * d));
  ld seg1 = 0.5L * r1*r1 * (alpha - sinl(alpha));
  ld seg2 = 0.5L * r2*r2 * (beta  - sinl(beta));
  return seg1 + seg2;
}
\end{minted}

\subsection{Line Sweep (Rectangle Union Area)}
Union area of axis-aligned rectangles via sweep line and segment tree. $O(N \log N)$.
\begin{minted}{c++}
struct Event {
  ll x, y1, y2, type;
  bool operator<(const Event &o) const { return x < o.x; }
};
struct SegTree {
  vector<int> cnt;
  vector<ll> len, ys;
  int n;
  SegTree(vector<ll> &ys_) : ys(ys_), n(ys_.size()-1) {
    cnt.assign(4*n, 0); len.assign(4*n, 0);
  }
  void upd(int i, int l, int r, int ql, int qr, int v) {
    if (ql >= r || qr <= l) return;
    if (ql <= l && r <= qr) cnt[i] += v;
    else {
      int m = (l+r)/2;
      upd(i*2, l, m, ql, qr, v);
      upd(i*2+1, m, r, ql, qr, v);
    }
    if (cnt[i]) len[i] = ys[r] - ys[l];
    else len[i] = (r-l==1) ? 0 : len[i*2] + len[i*2+1];
  }
  ll total() { return len[1]; }
};
ll union_area(vector<tuple<ll,ll,ll,ll>> &rects) {
  vector<Event> ev;
  vector<ll> ys;
  for (auto [x1,y1,x2,y2] : rects) {
    ev.push_back({x1, y1, y2, 1});
    ev.push_back({x2, y1, y2, -1});
    ys.push_back(y1); ys.push_back(y2);
  }
  sort(ev.begin(), ev.end());
  sort(ys.begin(), ys.end());
  ys.erase(unique(ys.begin(), ys.end()), ys.end());
  SegTree seg(ys);
  ll ans = 0, prev = ev[0].x;
  for (auto &e : ev) {
    ans += (e.x - prev) * seg.total();
    int l = lower_bound(ys.begin(),ys.end(),e.y1)-ys.begin();
    int r = lower_bound(ys.begin(),ys.end(),e.y2)-ys.begin();
    seg.upd(1, 0, seg.n, l, r, e.type);
    prev = e.x;
  }
  return ans;
}
\end{minted}

\subsection{Intersection of Many Cubes / 3D Boxes}
Given $n$ axis-aligned boxes $[x_{i,1}, x_{i,2}] \times [y_{i,1}, y_{i,2}] \times [z_{i,1}, z_{i,2}]$, find their mutual intersection volume.
The intersection along each axis is independent:
\[ L_x = \max_i x_{i,1},\quad R_x = \min_i x_{i,2},\quad \Delta x = \max(0\text{LL}, R_x - L_x) \]
Compute $\Delta y, \Delta z$ analogously. The intersection volume is:
\[ V = \Delta x \cdot \Delta y \cdot \Delta z \]
\textbf{Trap:} Always clamp each $\Delta \ge 0$ individually before multiplying. Never multiply $(R-L)$ directly, as two disjoint intervals with negative spans can falsely multiply to a positive volume.

